% Study sheet: Volume III, Chapter 1, Introduction.
% Numbers: \val{...} from introduction_studysheet_numbers.tex, written by introduction_studysheet_numbers.py.
\documentclass{studysheet}
\input{introduction_studysheet_numbers}
\sheetvolume{3}
\sheetchapter{1}{Introduction}
\sheetsources{Numbers from the chapter's mlsysim scenarios (hypothetical prices and step accuracies); questions marked assumed use new inputs.}

\begin{document}
% ===================================================================== FRONT: LEARN
\frontside
\begin{bigidea}
A foundation model only \hl{proposes}; a fluent proposal can be wrong. The \hl{runtime} around it
holds every permission, carries state across turns and checks the result. The unit of
design is the \hl{trajectory}; its horizon, state and authority size the runtime.
\end{bigidea}

\begin{leftcol}
\sheetsec{Key ideas}
\begin{ideas}
\idea{Who closes the loop.}{Predictive: code consumes a score. Chat: a person runs tools. Agent: the
  runtime dispatches only granted actions.}
\idea{The trajectory is the unit.}{The recorded states, proposals, decisions and observations of one task;
  cost and risk accrue over all of it.}
\idea{Four ways to end.}{Accepted, failed, unresolved (evidence short or budget spent) or escalated to
  a person. A silent crash is a defect.}
\idea{The fail-plausible fault.}{Output looks like success while the property fails; crash
  monitors see a normal return.}
\idea{Invariant closure.}{Bounds close mechanically below the model; acceptance needs independent
  evidence above it, against a contract fixed before the run.}
\idea{Six phases per turn; one is the model's.}{Assemble context, invoke model, validate and authorize,
  execute tool, record observation, assess against contract.}
\end{ideas}

\vspace{10pt}
\sheetsec{The H\textperiodcentered S\textperiodcentered A exposure space}
\begin{sheettable}{@{}c L c L@{}}
\multicolumn{2}{@{}l}{\thead{State carried}} & \multicolumn{2}{l@{}}{\thead{Authority}}\\ \midrule
$\StateExposureLevel_0$ & none: one call & $\AuthorityExposureLevel_0$ & read-only\\
$\StateExposureLevel_1$ & context, serving cache & $\AuthorityExposureLevel_1$ & edits a disposable workspace\\
$\StateExposureLevel_2$ & task-local working copy & $\AuthorityExposureLevel_2$ & external; undo ready before dispatch\\
$\StateExposureLevel_3$ & durable, shared, outlives the task & $\AuthorityExposureLevel_3$ & irreversible: approve before dispatch\\
\end{sheettable}
\vspace{3pt}
{\fontsize{8.6}{11}\selectfont\color{muted}Levels follow the environment,
not the command: one deletion is $\AuthorityExposureLevel_1$ in a sandbox, $\AuthorityExposureLevel_3$ on the only copy.\par}
\vspace{5pt}
\begin{sheettable}{@{}l L l L@{}}
\multicolumn{4}{@{}l@{}}{\thead{Lighthouses (turns, state, authority)}}\\ \midrule
Support & few, $\StateExposureLevel_1$, $\AuthorityExposureLevel_0$ & Remediation & tens, $\StateExposureLevel_3$, $\AuthorityExposureLevel_3$\\
Repair & hundreds, $\StateExposureLevel_2$, $\AuthorityExposureLevel_1$ & Research & thousands, $\StateExposureLevel_3$, $\AuthorityExposureLevel_0$\\
\end{sheettable}

\vspace{10pt}
\fcard{Five-part task contract}%
  {\TaskContract = \langle \TaskContractGoal,\ \TaskEnvironment,\ \PermittedActions,\ \AvailableObservations,\ \CompletionCriteria \rangle}%
  {Goal with a negative scope, environment, permitted actions, bounded observations, completion
   criteria. A prompt states intent; only the contract binds.}%
  {}{}
\end{leftcol}%
\begin{rightcol}
\sheetsec{Closing the loop from outside the model}
{\centering
\begin{tikzpicture}[sheet]
  \node[accbox, minimum width=236pt, minimum height=17pt] (up) at (118,96) {Above: accept only on evidence, against the contract};
  \node[hotbox, minimum width=168pt, minimum height=17pt] (m) at (86,70) {Model: proposes};
  \node[nbox, minimum width=168pt, minimum height=17pt] (dn) at (86,36) {Below: grant, sandbox, budgets};
  \node[nbox, dashed, minimum width=168pt, minimum height=15pt] (env) at (86,3) {Environment};
  \draw[arr] (40,61.5) -- (40,44.5) node[midway, right, note, xshift=2pt] {proposal};
  \draw[arr] (40,27.5) -- (40,10.5) node[midway, right, note, xshift=2pt] {allowed effect};
  \draw[arr] (170,3) -- (214,3) -- (214,87.5);
  \node[note, anchor=east, align=right] at (210,44) {observations,\\test results};
\end{tikzpicture}\par}
\vspace{5pt}
{\fontsize{8.6}{11}\selectfont Evidence, weakest first: self-report, static analysis, visible tests,
sealed tests (hidden, tamper-proof), formal proof. Each proves only what it measured; human review
covers the rest.\par}

\vspace{6pt}
\fcard{Horizon compounding}%
  {\Pr(\text{success}) = (1-\PerStepErrorRate)^{\TurnHorizon} \approx e^{-\PerStepErrorRate\TurnHorizon}}%
  {$\TurnHorizon$ unchecked independent steps, each wrong with probability $\PerStepErrorRate$.
   $\PerStepErrorRate\TurnHorizon$ sets reliability: doubling $\TurnHorizon$ squares success.}%
  {Checks with restoration shorten the unchecked horizon.}%
  {\begin{tikzpicture}[sheet]
     \foreach \yv in {0,50,100}{\draw[gridline] (30,{\yv*0.5}) -- (230,{\yv*0.5});
       \node[small, anchor=east] at (27,{\yv*0.5}) {\yv\%};}
     \foreach \h in {0,50,100,150,200}{\node[small, anchor=north] at ({30+\h},-3) {\h};}
     \draw[accent, line width=1.1pt] plot[domain=0:200, samples=80, smooth] ({30+\x},{50*exp(\val{lnstep}*\x)});
     \filldraw[fill=accentsoft, draw=accent] ({30+\val{hshort}},{0.5*\val{pshort}}) circle (2pt);
     \filldraw[fill=accent, draw=accent] ({30+\val{hlong}},{0.5*\val{plong}}) circle (2pt);
     \node[lab, anchor=south west] at ({34+\val{hshort}},{0.5*\val{pshort}+1}) {\val{hshort} turns: \val{pshort}\%};
     \node[lab, anchor=south west] at ({34+\val{hlong}},{0.5*\val{plong}+1}) {\val{hlong} turns: \val{plong}\%};
     \node[note, anchor=north] at (130,-12) {turns, at \val{stepok}\% correct per step};
   \end{tikzpicture}}
\fcard{Cost per accepted task}%
  {\CostPerAcceptedTask = \frac{\mathbb{E}[\AttemptCost]}{\AcceptanceProbability}, \qquad
   \text{pays if } \frac{\MechanismAttemptCost}{\MechanismAcceptanceProbability} < \frac{\AttemptCost}{\AcceptanceProbability}}%
  {Retries: $1/\AcceptanceProbability$ attempts expected, each charged in full. A mechanism, or a
   model-directed loop over a fixed workflow, pays only if it lowers the ratio under the same contract.}%
  {$\IndependentAttemptCount$ independent attempts succeed with $1-(1-\AcceptanceProbability)^{\IndependentAttemptCount}$;
   repeating one mistake gains nothing. Break it: fresh evidence, a new hypothesis, or a
   restart on a rolled-back workspace.}{}
\end{rightcol}

% ===================================================================== BACK: PRACTICE
\backside
\begin{answers}
\ans{Finish}{$\val{cmech} \div \val{pmechfrac}$ = \$\val{ceffmech}, above \$\val{ceff}: on
  cost alone, it loses. It pays only once $\MechanismAcceptanceProbability > \val{cmech} \div \val{ceff} \approx
  \val{pneedfrac}$ (cost card).}
\ans{1}{(a) $\val{qokfrac}^{\val{qh}} \approx \val{qpfrac}$: \val{qp}\%. (b) \val{qpfrac} squared, about \val{qptwo}\%:
  doubling the horizon squares success. (c) $\PerStepErrorRate\TurnHorizon \approx -\ln \val{qtargetfrac} =
  \val{qepsh}$, so $\PerStepErrorRate \approx \val{qeps}$\% per step (horizon card).}
\ans{2}{Lookup read-only; seat hold compensable, release ready before dispatch; refund irreversible.
  State: session context plus durable, shared booking records. Before the refund: approval below the
  model; a later check is too late (table).}
\columnbreak
\ans{3}{Self-report is the weakest level; the sealed test shows only that the single-threaded case passes,
  untampered. Require multi-threaded stress runs and review, per $\CompletionCriteria$; else
  unresolved, decided in phase six (closure).}
\ans{4}{(a) All fail with $\val{qmiss}^{\val{qk}} \approx \val{qallfail}$, so \val{qany}\% succeed. (b) Still \val{qpa}\%, and every failing task burns all
  \val{qk} attempts: the formula assumes independence. Break it: fresh evidence, a new hypothesis, or a
  clean restart with the workspace rolled back (cost card).}
\end{answers}

\begin{leftcol}
\sheetsec{Worked example}
\begin{practicebox}\boxtitle{Cheaper tokens or more acceptances?}
A repair attempt spends \$\val{cmodel} on tokens and \$\val{cother} on sandbox, tests and review;
\val{pacc}\% are accepted, retried until one is. Halve the token price, or raise
acceptance to \val{pbetter}\% with a tool fix?\par\vspace{5pt}
{\centering
\begin{tikzpicture}[sheet]
  \pgfmathsetmacro\sc{120/\val{ceff}}
  \foreach \lab/\v/\y/\sty in {baseline/\val{ceff}/34/nbox, half-price tokens/\val{ceffcheap}/17/nbox, more acceptance/\val{ceffbetter}/0/hotbox}{
    \node[lab, anchor=east] at (80,\y+6) {\lab};
    \draw[\sty] (84,\y) rectangle ({84+\v*\sc},\y+12);
    \node[note, anchor=west] at ({88+\v*\sc},\y+6) {\$\v};}
\end{tikzpicture}\par}\vspace{4pt}
\begin{steps}
\item Baseline: \val{cmodel} + \val{cother} = \$\val{cattempt} per attempt, $1/\val{paccfrac}$ = \val{attempts}
  attempts: \$\val{cattempt} $\div$ \val{paccfrac} = \res{\$\val{ceff}}.
\item Half-price tokens: \val{cheapmodel} + \val{cother} = \$\val{ccheap}; $\div$ \val{paccfrac} = \res{\$\val{ceffcheap}}.
\item More acceptance: \$\val{cattempt} $\div$ \val{pbetterfrac} = \res{\$\val{ceffbetter}}, narrowly best.
\item Break-even: full price matches \$\val{ceffcheap} at acceptance \val{cattempt} $\div$ \val{ceffcheap}
  $\approx$ \res{\val{pbreakeven}\%}.
\end{steps}
\takeaway{A price cut removes only its own share; acceptance divides the whole attempt.}
\end{practicebox}

\vspace{10pt}
\sheetsec{Your turn: finish it}
\begin{practicebox}\boxtitle{Does a sealed-test stage pay for itself?}
\assumed{} same repair; a sealed test stage adds \$\val{mechadd} per attempt and raises acceptance from
\val{pacc}\% to \val{pmech}\%.\par\vspace{4pt}
\textbf{Given:} new attempt \val{cattempt} + \val{mechadd} = \$\val{cmech}.\par
\textbf{You:} find the new cost per accepted task, and the acceptance the stage needs to beat \$\val{ceff}.\par\vspace{30pt}
\end{practicebox}
\end{leftcol}%
\begin{rightcol}
\sheetsec{Common mistakes}
\mistake{``The agent says the bug is fixed and the tests pass, so it is done.''}%
  {A fail-plausible run writes the same report. At \val{stepok}\% per step, an unchecked \val{hlong}-turn repair
   is right in \val{plong}\% of runs; if every run reports success, about \val{wrongshare}\% of acceptances are wrong.}{}
\mistake{``Any variation in the environment needs a model-directed loop.''}%
  {Cloning, installing and linting never change order: a fixed workflow is cheaper. A loop earns its
   tokens only where it lowers $\CostPerAcceptedTask$, such as an unknown race.}{}
\mistake{``A large context window removes the need to manage state.''}%
  {Over \val{ctxturns} turns (\val{ctxfirst} tokens, plus \val{ctxper} per turn) the last call reads
   \val{ctxlast}, \val{ctxshare}\% of a \val{ctxwindow} window, yet the calls read \val{ctxcum} tokens in all
   (\val{ctxreread}$\times$), stale file versions included.}{}

\vspace{2pt}
\sheetsec{Check yourself}
\begin{qs}
\item \assumed{} \val{qok}\% of steps correct, independent, unchecked. (a) Task
  success at \val{qh} turns? (b) At \val{qhtwo}? (c) What per-step error holds \val{qtarget}\% at \val{qh} turns?
\item The support agent can look up a booking, hold a seat (released if unused) and refund to the
  customer's bank. Place each action on the authority axis. Which state levels? What must
  the runtime add before the refund?
\item A repair agent ends with ``race condition fixed; all tests pass.'' The log shows one sealed,
  single-threaded test exited zero. What does this support? What would you require, and in which phase?
\item \assumed{} each attempt is accepted with probability \val{qpa}\%. (a) Chance \val{qk}
  independent attempts succeed? (b) Each failure is appended to the prompt and the model repeats its mistake.
  What is the success chance now? What breaks the repetition?
\end{qs}
\end{rightcol}
\end{document}
